% PLOS Computational Biology manuscript
% Uses standard article class with PLOS formatting requirements:
% double-spaced, line numbers, single column, Vancouver references

\documentclass[10pt,letterpaper]{article}

% --- Required packages ---
\usepackage[utf8]{inputenc}
\usepackage[T1]{fontenc}
\usepackage{amsmath,amssymb}
\usepackage{graphicx}
\usepackage{booktabs}
\usepackage{array}
\usepackage{url}
\usepackage[hidelinks]{hyperref}
\usepackage{lineno}       % Line numbers (PLOS requirement)
\usepackage{setspace}     % Double spacing (PLOS requirement)
\usepackage[margin=1in]{geometry}
\usepackage[numbers,sort&compress]{natbib}  % Vancouver-style numbered references

% --- PLOS formatting ---
\doublespacing
\linenumbers

% Vancouver style: numbered references in square brackets, ordered by appearance
\bibliographystyle{unsrtnat}

% No headers/footers except page numbers
\pagestyle{plain}

\begin{document}

% ===========================
% TITLE PAGE
% ===========================

\begin{flushleft}
{\Large\bfseries Causal intervention validation of gene regulatory signals in scGPT}
\bigskip

Ihor Kendiukhov$^{1,*}$

\bigskip

\textbf{1} Department of Computer Science, University of Tuebingen, Sand~14,
72076 Tuebingen, Germany

\bigskip

$*$ Corresponding author

E-mail: kenduhov.ig@gmail.com

\end{flushleft}

\newpage

% ===========================
% ABSTRACT
% ===========================

\section*{Abstract}

We evaluate whether causal interventions on scGPT gene tokens recover known
transcription factor (TF)--target dependencies beyond attention-only heuristics.
Using Tabula Sapiens kidney, lung, and immune subsets plus an external lung
atlas, we ablate or swap TF token values and quantify changes in target-token
readouts. We compare intervention scores to curated gene regulatory references
(TRRUST, DoRothEA), include attention and coexpression baselines, trace
component-level circuits via activation patching, and validate on CRISPR
perturbation data from Adamson.
Kidney multi-seed runs show modest enrichment (AUPR~$\approx$~0.60), while
lung results are substantially stronger (AUPR up to 0.78, permutation
$p{=}0.001$) and remain positive in higher-pair settings. Mechanistic tracing
localizes effects to a small set of attention heads and mid-to-late MLP blocks.
External perturbation validation remains low power but provides an independent
check of causal signals. We present the full pipeline, quantitative evidence,
and an honest assessment of current limitations. Source code and evaluation
artifacts are available at
\url{https://github.com/Biodyn-AI/casual-intervention} (to be made public upon acceptance).

% ===========================
% AUTHOR SUMMARY (PLOS requirement)
% ===========================

\section*{Author summary}

We wanted to understand whether a large AI model trained on gene expression
data from millions of human cells has actually learned how genes regulate each
other, or whether it has only memorized statistical patterns. To test this, we
deliberately disrupted one gene's representation inside the model and measured
how much this changed the model's predictions for known regulatory targets.
We found that the answer depends on the tissue. In lung tissue, the model
appears to have learned genuine regulatory relationships that go beyond simple
correlations, with disruptions to transcription factor genes causing
significant changes in their known targets. In kidney and immune tissues, the
signal was weaker. We also traced which internal components carry these
regulatory signals, finding that a small number of attention heads and neural
network layers are responsible. Our work provides a practical toolkit for
auditing whether biological AI models have learned real biology or just
correlations, which matters as these models are increasingly used for drug
target discovery and cell type classification.

% ===========================
% INTRODUCTION
% ===========================

\section*{Introduction}

Transformer-based \cite{vaswani_attention_2017} single-cell foundation models,
including scGPT \cite{cui_scgpt_2023}, Geneformer
\cite{theodoris_geneformer_2023}, scBERT \cite{yang_scbert_2022}, and
scFoundation \cite{hao_scfoundation_2024}, capture gene dependencies at
scale, but attention weights alone are correlational and do not establish
causal influence in the sense of Pearl's interventionist framework
\cite{pearl_causality_2009}. Prior gene regulatory network (GRN) inference
methods such as SCENIC \cite{aibar_scenic_2017} and the Inferelator
\cite{bonneau_grn_review_2007} derive regulatory structure from expression
data using motif enrichment and coexpression, and systematic GRN benchmarks
\cite{pratapa_beeline_2020} have evaluated these methods, yet they do not test
whether a model's internal representations encode directional TF--target
relationships. We ask whether direct causal interventions on scGPT token
representations produce scores aligned with curated TF--target edges, and
whether mechanistic tracing can localize the internal components responsible
for those effects.

Our approach draws on causal mediation analysis
\cite{vig_causal_mediation_2020} and activation patching techniques developed
for language model interpretability
\cite{meng_rome_2022,goldowsky_dill_path_patching_2023}, adapting them to the
biological token structure of single-cell transformers. The mechanistic
interpretability paradigm \cite{olah_zoom_in_2020} seeks to decompose neural
network computations into understandable circuits, and recent automated methods
\cite{conmy_acdc_2023} have begun to scale this approach. Superposition
\cite{elhage_superposition_2022} poses a challenge for such analyses, as
features may share dimensions, but the relatively sparse gene-token structure
of scGPT may partially mitigate this concern. Related formalisms include causal
scrubbing \cite{chan_causal_scrubbing_2022}, which provides a complementary
framework for testing circuit-level hypotheses.

Our evaluation focuses on two complementary signals: curated regulatory
interactions (TRRUST and DoRothEA) and perturbation-derived edges from CRISPR
screening data \cite{dixit_perturbseq_2016,datlinger_perturbseq_2017}. The
growing availability of large-scale perturbation atlases
\cite{replogle_perturbseq_2022,norman_crispr_2019} provides increasingly rich
ground truth for such evaluations. The first signal provides a broad reference
derived from literature curation, while the second provides experimentally
induced causal signals. By aligning causal intervention scores with both, we
probe whether internal gene-token dependencies show evidence of causal
structure beyond correlation.

We introduce a causal intervention pipeline for scGPT that provides
cross-tissue and multi-seed results, compares against attention and
coexpression baselines, localizes circuitry with activation patching, and
reports a perturbation-based validation using Adamson CRISPR data.

% ===========================
% RESULTS (before Methods per PLOS convention)
% ===========================

\section*{Results}

\subsection*{Kidney multi-seed results}

Across four kidney seeds, TRRUST ablation achieves mean AUPR 0.6018 with
bootstrap 95\% CI [0.575, 0.628], while DoRothEA ablation reaches 0.5369
[0.505, 0.563]. Swap interventions are weaker (TRRUST 0.537, DoRothEA 0.509).
Evidence coverage is consistent (165--175 pairs from 240 sampled), stabilizing
estimates. The seed-level spread is modest (0.57--0.63 for TRRUST ablation),
indicating reproducibility at current scale (Fig~\ref{fig:kidney-seed-ci}).

The fact that ablation consistently outperforms swap deserves attention: zeroing
out a TF token value removes that gene's contribution entirely, producing a
clean counterfactual, whereas swapping replaces it with another gene's value,
which may partially preserve the original regulatory signal if the replacement
gene participates in similar pathways. In other words, ablation tests
\textit{necessity} of the TF representation, while swap is a noisier probe that
conflates removal with substitution. The gap between the two suggests that
scGPT encodes at least some TF-specific information that cannot be replaced by
an arbitrary gene token, which is a prerequisite for genuine causal structure
rather than generic positional or cell-state effects.

% Fig 1 caption
\begin{figure}[!ht]
\caption{{\bf TRRUST ablation AUPR by kidney seed with bootstrap 95\% CIs.}
Each point represents one random seed; error bars are bootstrap 95\% confidence
intervals. The dashed line indicates the random baseline (AUPR~=~0.50).}
\label{fig:kidney-seed-ci}
\end{figure}

The modest absolute AUPR values (around 0.60 against a random baseline of
$\sim$0.50 in balanced settings) reflect the inherent difficulty of
recovering directed regulatory edges from a model that was trained on
expression reconstruction, not on regulatory inference per se. Even a
partially enriched signal is informative because it suggests that the model's
internal representations are not agnostic to regulatory relationships.

\subsection*{Cross-tissue performance}

Lung shows substantially stronger enrichment than kidney or immune
(Table~\ref{tab:cross-tissue}). Lung TRRUST ablation AUPR reaches 0.7272
(permutation $p{=}0.001$) and DoRothEA 0.7820 ($p{=}0.001$). Immune results are
weaker and nonsignificant, while the Krasnow lung dataset shows mixed
alignment consistent with dataset shift. The high-pair lung run (160+160)
retains strong performance (TRRUST 0.7574, $p{=}0.004$), indicating robustness
to expanded candidate sets (Fig~\ref{fig:cross-tissue}).

The tissue-level variation has a plausible biological explanation. Lung tissue
in Tabula Sapiens contains well-characterized epithelial and endothelial cell
types with active, tissue-specific transcription factor programs (e.g., NKX2-1,
FOXJ1, SOX family members) that are richly represented in curated databases.
Kidney has a more heterogeneous cell composition with many transporter and
metabolic genes that may be less well-covered by TF-centric references like
TRRUST. Immune cells, while transcriptionally active, often rely on
combinatorial and context-dependent regulation (e.g., cytokine signaling
cascades) that single-TF ablation may not capture cleanly.

The Krasnow lung dataset, collected with Smart-seq2 rather than 10x Chromium,
shows lower TRRUST alignment (AUPR 0.36) despite being lung tissue. This
likely reflects platform effects on the expression distributions that scGPT was
pretrained on: the model's whole-human checkpoint was trained predominantly on
droplet-based data, and the distributional shift from Smart-seq2 full-length
transcripts may degrade the learned TF representations. The fact that DoRothEA
alignment remains moderately positive (0.63) on the same dataset suggests that
some regulatory signals survive the platform shift, particularly the
higher-confidence interactions captured by DoRothEA's literature scoring.

\begin{table}[!ht]
\centering
\caption{{\bf Cross-tissue ablation results and robustness runs.}
Bold $p$-values indicate significance at $\alpha{=}0.05$.}
\label{tab:cross-tissue}
\begin{tabular}{llcc}
\toprule
\textbf{Group} & \textbf{Ref} & \textbf{AUPR} & \textbf{Perm $p$} \\
\midrule
Lung & TRRUST & 0.7272 & \textbf{0.001} \\
Lung & DoRothEA & 0.7820 & \textbf{0.001} \\
Immune & TRRUST & 0.5955 & 0.359 \\
Immune & DoRothEA & 0.5267 & 0.672 \\
Krasnow lung & TRRUST & 0.3635 & 0.934 \\
Krasnow lung & DoRothEA & 0.6255 & 0.769 \\
Lung (fast) & TRRUST & 0.6543 & 0.369 \\
Lung (fast) & DoRothEA & 0.6915 & 0.209 \\
Lung (high-pair) & TRRUST & 0.7574 & \textbf{0.004} \\
Lung (high-pair) & DoRothEA & 0.6910 & 0.102 \\
\bottomrule
\end{tabular}
\end{table}

% Fig 2 caption
\begin{figure}[!ht]
\caption{{\bf Cross-tissue ablation AUPR across tissues and robustness runs.}
Bar heights show AUPR for ablation intervention; error bars indicate bootstrap
95\% CIs where multiple seeds are available.}
\label{fig:cross-tissue}
\end{figure}

\subsection*{Baselines versus causal scores}

Causal scores outperform attention and coexpression baselines for lung and
high-pair lung runs, while kidney and immune are competitive with baselines
(Table~\ref{tab:baseline}). The lung improvement is substantial: causal AUPR
exceeds attention by +0.184 and coexpression by +0.140. Kidney gains are
small (+0.017), and immune shows negative deltas. Reduced-gene runs reinforce
this: reduced lung remains above attention (0.674 vs.\ 0.608, $p{=}0.007$),
while reduced kidney trails attention (Fig~\ref{fig:baseline}).

The relationship between these three scoring methods reveals something
important about what the model has learned. Attention scores capture which
tokens the model ``looks at'' during forward computation, but high attention
does not imply causal influence because attention heads can attend to tokens
for context without propagating regulatory information. Coexpression
correlations capture statistical co-occurrence of gene pairs across cells,
which can arise from shared upstream regulation, common cell-state programs,
or true direct regulation. Causal intervention scores, by contrast, test
whether \textit{removing} a specific TF token changes the model's prediction
at the target position, which is a more direct test of functional dependency
within the model's learned representation.

The fact that causal scores exceed both baselines in lung but not in kidney or
immune suggests that the model has learned genuinely directional TF--target
relationships for lung-specific programs, whereas in kidney and immune tissues,
the model's gene representations may be dominated by broader cell-state
signatures that correlational methods capture equally well. This is consistent
with the intuition that causal interventions are most informative when the
underlying signal has clear directionality, rather than being a diffuse
cell-type marker.

\begin{table}[!ht]
\centering
\caption{{\bf TRRUST ablation: causal scores vs.\ attention and coexpression baselines.}
AUPR values shown. Bold $p$-values indicate significance at $\alpha{=}0.05$.}
\label{tab:baseline}
\begin{tabular}{lcccc}
\toprule
\textbf{Group} & \textbf{Causal} & \textbf{Attn} & \textbf{Coex} & \textbf{$p$} \\
\midrule
Kidney & 0.602 & 0.585 & 0.583 & 0.105 \\
Lung & \textbf{0.727} & 0.543 & 0.587 & \textbf{0.001} \\
Immune & 0.596 & 0.617 & 0.628 & 0.359 \\
Kidney (red.) & 0.540 & 0.613 & 0.565 & 0.264 \\
Lung (red.) & 0.674 & 0.608 & 0.615 & \textbf{0.007} \\
Lung (hi-pair) & \textbf{0.757} & 0.589 & 0.729 & \textbf{0.004} \\
\bottomrule
\end{tabular}
\end{table}

% Fig 3 caption
\begin{figure}[!ht]
\caption{{\bf TRRUST ablation: causal scores vs.\ baselines across tissues.}
Grouped bar chart comparing causal intervention AUPR against attention and
coexpression baselines for each tissue and robustness configuration.}
\label{fig:baseline}
\end{figure}

\subsection*{Mechanistic case studies}

Activation patching localizes specific TF--target effects to a small set of
attention heads and MLP blocks. Example pairs
(GATA3${\rightarrow}$GATA3, FOS${\rightarrow}$FOS,
JUN${\rightarrow}$JUN, JUN${\rightarrow}$TNF) consistently highlight
mid-to-late MLPs and a handful of heads (Fig~\ref{fig:case-study},
Table~\ref{tab:case-studies}). Self-targeting pairs include late MLP blocks
(L8--L11), whereas cross-gene JUN${\rightarrow}$TNF concentrates on earlier
blocks (L0--L5), suggesting different interaction types engage distinct depth
regimes. Head lists are dominated by H0/H1 across layers L0--L7, indicating a
small subset carries most restoration signal.

The depth separation between self-targeting and cross-gene interactions aligns
with a natural hypothesis about how the model processes gene-gene
relationships. Self-targeting effects (a TF influencing its own output token)
can be resolved through relatively local, residual-stream computations that
accumulate through later MLP blocks, since the source and target share the same
token position. Cross-gene effects like JUN${\rightarrow}$TNF require
inter-token communication through attention, which must happen earlier in the
network so that downstream MLP blocks can integrate the transferred
information. The observation that cross-gene pairs rely on earlier components
is consistent with transformer architectures where attention moves information
between positions and MLPs transform it within positions.

These circuit-level findings, while based on a small set of case-study pairs,
provide a template for future mechanistic audits. If specific heads or MLP
blocks consistently mediate known regulatory effects, they could serve as
targets for surgical interventions to test whether the model's regulatory
knowledge is localized or distributed, a question with practical implications
for model editing and fine-tuning strategies in biological foundation models.

% Fig 4 caption
\begin{figure}[!ht]
\caption{{\bf Case-study restoration strength by component type.}
Bars show mean restoration fraction for attention heads vs.\ MLP blocks across
case-study TF--target pairs. Higher values indicate greater contribution to
recovering the clean-run target readout.}
\label{fig:case-study}
\end{figure}

\begin{table}[!ht]
\centering
\caption{{\bf Top contributing heads and MLP blocks from causal tracing.}}
\label{tab:case-studies}
\begin{tabular}{lp{4cm}p{5cm}}
\toprule
\textbf{Pair} & \textbf{Top heads} & \textbf{Top MLP blocks} \\
\midrule
GATA3${\to}$GATA3 & L6-H1, L1-H1, L2-H1 & L9, L8, L10, L4 \\
FOS${\to}$FOS & L0-H0, L5-H1, L1-H0 & L10, L1, L8, L9 \\
JUN${\to}$JUN & L7-H0, L2-H1, L4-H1 & L0, L11, L9, L6 \\
JUN${\to}$TNF & L3-H0, L1-H0, L2-H1 & L3, L0, L2, L4 \\
\bottomrule
\end{tabular}
\end{table}

\subsection*{Perturbation validation}

Adamson perturbation validation remains low power at current scale. The scaled
run (40+40 pairs) yields AUPR 0.435 for ablation and 0.502 for swap; the
large run (100+100 pairs) yields 0.498 and 0.500. None are significant under
permutation tests (Table~\ref{tab:perturb}). As pair counts increase, AUPR
drifts toward chance, which we believe reflects several compounding factors
rather than a single failure mode.

First, Adamson perturbation-derived edges are defined by the top-$k$ expression
changes after CRISPR knockout, which captures the strongest downstream
\textit{transcriptional} effects. However, scGPT's causal intervention operates
at the \textit{representational} level, testing whether removing a gene token's
value changes the model's output at another token position. These are
conceptually related but not identical: a gene that is a strong transcriptional
target of a TF in real cells may not be the gene whose model representation is
most disrupted by ablating the TF token, especially if the model has learned to
represent the target through cell-state features rather than direct
TF-mediated effects.

Second, the Adamson dataset uses a CRISPRi screen in K562 cells (a chronic
myelogenous leukemia line), which is biologically distant from the primary
tissues on which scGPT was pretrained. The model may encode regulatory
relationships specific to normal tissue contexts that do not align with
K562-specific perturbation effects.

Third, the evidence filter (requiring both TF and target to co-occur in the
scGPT evaluation cells) substantially reduces the effective evaluation set,
compressing statistical power. At 75--175 evidence pairs with near-chance AUPR,
the permutation test lacks the resolution to detect a weak true signal even if
one exists.

\begin{table}[!ht]
\centering
\caption{{\bf Perturbation validation (Adamson).}
Int.\ = intervention type, $n$ = evidence pairs.}
\label{tab:perturb}
\begin{tabular}{llcccc}
\toprule
\textbf{Run} & \textbf{Int.} & \textbf{AUPR} & \textbf{AUROC} & \textbf{$p$} & \textbf{$n$} \\
\midrule
scaled & abl & 0.435 & 0.416 & 0.910 & 75 \\
scaled & swap & 0.502 & 0.503 & 0.479 & 75 \\
large & abl & 0.498 & 0.491 & 0.606 & 175 \\
large & swap & 0.500 & 0.513 & 0.589 & 175 \\
pilot & abl & 0.576 & 0.483 & 0.430 & 31 \\
pilot & swap & 0.577 & 0.521 & 0.402 & 31 \\
\bottomrule
\end{tabular}
\end{table}

\subsection*{High-confidence DoRothEA (A/B)}

The stricter DoRothEA A/B subset reduces label breadth, depressing absolute
AUPR, but lung retains a clear permutation signal: kidney mean AUPR 0.396
(median $p{=}0.085$) versus lung AUPR 0.482 ($p{=}0.001$), consistent with the
tissue-specific differences observed throughout. The DoRothEA A/B confidence
levels correspond to interactions supported by direct experimental evidence
(chromatin immunoprecipitation, reporter assays) rather than computational
predictions, so the persistent lung signal under this stricter filter
strengthens the case that the model captures bona fide regulatory
relationships rather than spurious correlations with loosely curated edges.

% ===========================
% DISCUSSION
% ===========================

\section*{Discussion}

Our results paint a nuanced picture of what single-cell foundation models
learn about gene regulation. Causal ablations yield modest enrichment on kidney
and substantially stronger performance on lung, with significant TRRUST
enrichment across multiple lung configurations. This pattern is consistent
across reference databases, robustness runs, and high-confidence subsets,
which makes it unlikely to be an artifact of a single evaluation choice.

The tissue dependence of causal signal strength carries an important
implication for how we think about foundation model representations. Rather
than encoding a universal regulatory program, scGPT appears to have learned
tissue-specific regulatory structure that is stronger where training data
provided richer and more coherent transcriptional programs. Lung, with its
well-characterized epithelial differentiation hierarchies, may offer cleaner
TF-target signals for the model to internalize during pretraining, whereas
kidney and immune tissues present more complex or combinatorial regulatory
logic that a token-level ablation may not isolate cleanly.

The baseline comparisons add an important calibration. In tissues where causal
scores do not clearly exceed attention or coexpression, the intervention
pipeline is not providing additional information beyond what simple
correlational proxies capture. This does not invalidate the approach, but it
does suggest that causal interventions are most valuable in a targeted
setting: when there is reason to believe a specific tissue or TF program is
well-represented in the model, intervention-based scoring can extract
directional information that correlational methods miss.

Component-level tracing provides interpretable circuits for selected
TF--target pairs. The concentration of restoration effects in a small set of
heads and MLPs suggests a tractable mechanistic substrate within scGPT,
offering targets for focused audits and hypothesis-driven validation. The
depth separation between self-targeting and cross-gene interactions hints at
a functional organization within the transformer that parallels the
distinction between local autoregulatory feedback and distal regulatory
influence.

Perturbation validation results remain inconclusive, which we attribute
primarily to the biological distance between the K562 CRISPR screen and the
normal-tissue contexts on which the model was evaluated, compounded by limited
statistical power from evidence filtering. Future work should prioritize
perturbation datasets from primary tissues matching the evaluation context,
such as Perturb-seq in lung epithelial cells, to provide a fairer test of
whether the model's causal signals align with experimentally induced effects.

\paragraph{Limitations.}
Our evaluation has several limitations that scope the conclusions. All runs
use CPU-limited subsets (30--120 cells per run) with moderate pair counts
(30--160 positive pairs), which limits statistical power particularly for the
perturbation validation where only 75--175 evidence pairs are available. The
readout metric is a single target-token scalar (mean across hidden features),
which may not fully capture the dimensionality of regulatory effects encoded
in the model's representations. The perturbation validation relies on the
Adamson dataset, which uses a CRISPRi screen in K562 cells that are
biologically distant from the primary tissues used in the main evaluation;
this mismatch likely contributes to the inconclusive perturbation results.
Gene symbol mapping from Ensembl identifiers to HGNC symbols can introduce
alias errors, potentially misaligning a small fraction of evaluation pairs.
The mechanistic case studies cover only four TF--target pairs and may not
generalize to the full range of regulatory interactions. Finally, our
evaluation uses a single foundation model (scGPT); extending to other
architectures is needed to determine whether the observed patterns are
model-specific or reflect general properties of expression-trained
transformers.

% ===========================
% MATERIALS AND METHODS (after Results/Discussion per PLOS convention)
% ===========================

\section*{Materials and methods}

\subsection*{Model and datasets}

We use the scGPT whole-human checkpoint from the model zoo, loaded with
standard model arguments. Fast transformer backends are disabled for
deterministic patching. All runs are CPU-only with fixed random seeds.

We use Tabula Sapiens kidney, lung, and immune subsets
\cite{tabula_sapiens_2022}, an external lung atlas from
Travaglini et al.\ \cite{travaglini_lung_2020}, and the Adamson CRISPR
perturbation dataset \cite{adamson_dataverse_2022}. All datasets were
originally generated using droplet-based protocols
\cite{zheng_10x_2017} except the Krasnow lung atlas (Smart-seq2).
Preprocessing follows standard Scanpy \cite{wolf_scanpy_2018} workflows
(normalize\_total=10\,000, log1p), selects 5\,000 highly variable genes
restricted to the scGPT vocabulary, and normalizes gene symbols via HGNC alias
mappings. We note that deep learning denoising approaches
\cite{eraslan_sc_review_2019} and integration methods
\cite{stuart_seurat_2019} were not applied, as we sought to evaluate the
foundation model's representations on minimally processed data.

\subsection*{Intervention pairs and scoring}

Positive labels come from curated TF--target edges in TRRUST \cite{han_trrust_2015}
and DoRothEA \cite{dorothea_github}; negatives are random controls from the same gene sets.
For each TF--target pair and each cell containing both genes, we compute a baseline readout and an
intervened readout at the target token position, reduced by mean across
features. We apply value ablation (TF token value zeroed) and token-value swap
interventions. The pair score is the mean absolute effect
$|\bar{\Delta}| = |y_{\text{baseline}} - y_{\text{intervened}}|$ across cells.

\subsection*{Causal tracing}

We patch clean activations into ablated forward passes and measure restoration
of the target readout, yielding per-layer, per-head, and per-MLP restoration
scores for selected case-study pairs.

\subsection*{Baselines and metrics}

We compare causal scores against (i)~attention scores aggregated across all
layers and heads and (ii)~coexpression correlations computed from normalized
expression. We compute AUPR and AUROC \cite{davis_prauc_2006} on pairs with
evidence; significance is
estimated by 1\,000 permutations of the labels, reporting the fraction of
permutations with AUPR $\geq$ observed (median across seeds where applicable).
We additionally report DoRothEA A/B subsets for high-confidence evaluation.

\subsection*{Perturbation validation}

We derive perturbation edges from Adamson \cite{adamson_cell_2016,adamson_dataverse_2022}
by selecting the top-$k$ ($k{=}50$) absolute expression deltas versus control per perturbation.
We evaluate causal scores on balanced pair sets: 40+40 (scaled) and 100+100 (large), with
max\_cells=60 per run.

\subsection*{Experimental configurations}

We evaluate kidney (four seeds), lung, and immune subsets with 120 max cells
and 120+120 candidate pairs. Robustness runs include a high-pair lung setting
(160+160 pairs), reduced-gene runs, and a Krasnow lung dataset. Adamson
experiments use scaled (40+40) and large (100+100) pair counts.
Table~\ref{tab:run-configs} summarizes configurations.

\begin{table}[!ht]
\centering
\caption{{\bf Run configurations.}
TS = Tabula Sapiens. Pairs with evidence are a subset filtered to cells
containing both source and target.}
\label{tab:run-configs}
\begin{tabular}{llcc}
\toprule
\textbf{Run group} & \textbf{Dataset} & \textbf{Cells} & \textbf{Pairs} \\
\midrule
Kidney (4 seeds) & TS kidney & 120 & 120+120 \\
Lung & TS lung & 120 & 120+120 \\
Immune (2 seeds) & TS immune & 120 & 120+120 \\
Krasnow lung & Krasnow & 40 & 30+30 \\
Lung (fast) & TS lung & 60 & 60+60 \\
Lung (high-pair) & TS lung & 40 & 160+160 \\
Kidney (reduced) & TS kidney & 300 & 80+80 \\
Lung (reduced) & TS lung & 240 & 80+80 \\
Adamson (scaled) & Adamson & 60 & 40+40 \\
Adamson (large) & Adamson & 60 & 100+100 \\
\bottomrule
\end{tabular}
\end{table}

% ===========================
% CONCLUSION
% ===========================

\section*{Conclusion}

Causal interventions on scGPT token representations enrich for curated
TF--target edges, particularly in lung, and mechanistic tracing localizes
components for specific pairs. The strongest result is lung tissue AUPR up to
0.78 against DoRothEA references with permutation $p{=}0.001$, a signal
that survives across pair counts, robustness runs, and high-confidence
reference subsets. Kidney enrichment is modest but reproducible across seeds,
and immune tissue remains at baseline, together delineating where the model
has and has not learned directional regulatory structure.

Several directions follow naturally from these findings. First, scaling runs
to GPU-enabled full-tissue evaluations would increase pair counts by an order
of magnitude, sharpening the statistical resolution for tissues where current
enrichment is marginal. Second, replacing the mean-across-features readout
with richer target representations, such as per-feature cosine similarity or
learned probes, could recover regulatory effects that a single scalar
compresses away. Third, applying the same causal intervention framework to
other single-cell foundation models (Geneformer, scFoundation) would reveal
whether tissue-dependent regulatory encoding is architecture-specific or a
general property of expression-trained transformers. Finally, perturbation
validation using tissue-matched Perturb-seq datasets, particularly from lung
epithelial cells where model performance is strongest, would close the gap
between representational causal effects and experimental ground truth that the
Adamson K562 comparison currently leaves open.

More broadly, the pipeline we describe, combining token-level causal
interventions with component-level activation patching, provides a general
template for mechanistic audits of biological foundation models. As these
models are increasingly deployed for cell-type annotation, perturbation
prediction, and drug target prioritization, understanding which internal
components encode which biological relationships becomes a practical safety
and reliability concern, not merely an academic curiosity. The artifacts and
evaluation framework accompanying this work are designed to lower the barrier
for such audits across tissues, models, and regulatory contexts.

% ===========================
% ACKNOWLEDGMENTS
% ===========================

\section*{Acknowledgments}

This work was conducted independently without external funding.

% ===========================
% DATA AND CODE AVAILABILITY
% ===========================

\section*{Data availability statement}

All primary datasets used in this study are publicly available. The Tabula
Sapiens atlas \cite{tabula_sapiens_2022} is available at
\url{https://tabula-sapiens-portal.ds.czbiohub.org}. The Krasnow lung atlas
\cite{travaglini_lung_2020} is available through the Human Lung Cell Atlas at
\url{https://hlca.ds.czbiohub.org}. The Adamson CRISPR perturbation dataset
\cite{adamson_cell_2016,adamson_dataverse_2022} is available on Harvard
Dataverse at \url{https://doi.org/10.7910/DVN/BIKGRL}. TRRUST v2
\cite{han_trrust_2015} is available at
\url{https://www.grnpedia.org/trrust}. DoRothEA regulons
\cite{dorothea_github} are available through the decoupler Python package at
\url{https://github.com/saezlab/decoupler-py}. The scGPT whole-human
pretrained checkpoint \cite{cui_scgpt_2023} is available at
\url{https://github.com/bowang-lab/scGPT}.

All source code for the causal intervention pipeline, evaluation scripts,
run configurations, and figure generation code is available at
\url{https://github.com/Biodyn-AI/casual-intervention} (to be made public upon
acceptance). Causal scores, intermediate outputs, and reproducibility
artifacts are included in the repository.

% ===========================
% REFERENCES
% ===========================

\bibliography{references}

% ===========================
% SUPPORTING INFORMATION
% ===========================

\section*{Supporting information}

\paragraph*{S1 Table.}
\label{S1_Table}
{\bf Full causal metrics summary across all run configurations.}
Complete AUPR, AUROC, and permutation $p$-values for every run group,
reference database, and intervention type.

\paragraph*{S1 Data.}
\label{S1_Data}
{\bf Source code and evaluation artifacts.}
Available at \url{https://github.com/Biodyn-AI/casual-intervention} (to be made
public upon acceptance). Includes run configurations, evaluation scripts,
causal scores, and figure generation code.

\end{document}
